\begin{lemma}
Let $U\subseteq\mathbb{C}$ be open, non-empty and connected, let $m\ge0$ be an
integer, and
let $g\colon\mathbb{C}\to\mathbb{C}$ be $C^\infty$ on $U$.  Assume that
\[
  D^\nu g(z)=0\qquad\text{for every }z\in U\text{ and every }\nu\text{ with }
  |\nu|=m+1,
\]
where $D^\nu$ runs over the iterated real partial derivatives
$\partial_{j_1}\cdots\partial_{j_{m+1}}$ of \noderef{PartialWord},
$\partial_0=\partial/\partial x$ and $\partial_1=\partial/\partial y$.  Then
there are coefficients $a_{ij}\in\mathbb{C}$, $i,j\ge0$, with $a_{ij}=0$
whenever $i+j>m$, such that
\[
  g(z)=\sum_{i=0}^{m}\sum_{j=0}^{m}a_{ij}\,z^{i}\bar z^{\,j}
  \qquad(z\in U).
\]

This is the step Section 7 of the paper states as ``in each case, we conclude
that $f$ must be a homogeneous polynomial''.  Only the polynomial character is
proved here; the homogeneity, which turns the degree bound into an exact
degree, comes from the homogeneity (57) of the extremal inverse and is applied
separately.
\end{lemma}
\begin{proof}
Throughout, $e_0=1$ and $e_1=i$ are the two coordinate directions, so that
$\partial_j h=\mathrm{d}h[e_j]$, and for $h\in\mathbb{C}$ we write
$h^{(0)}=\operatorname{Re}h$ and $h^{(1)}=\operatorname{Im}h$, so that
$h=h^{(0)}e_0+h^{(1)}e_1$.

\emph{Step 1: the directional-derivative expansion.}  Let $\Omega\subseteq U$ be
open, let $h\in\mathbb{C}$, let $w\in\Omega$, and suppose the segment
$\{w+th:0\le t\le1\}$ lies in $\Omega$.  Put $\psi(t)=g(w+th)$.  We claim that
for every $s\ge0$ and every $t$ with $w+th\in\Omega$,
\begin{equation}\tag{1}
  \psi^{(s)}(t)=\sum_{|\nu|=s}c_\nu(h)\,D^\nu g(w+th),
\end{equation}
where the sum is over the $2^s$ words $\nu=(j_1,\dots,j_s)$ of length $s$ in
$\{0,1\}$ and $c_\nu(h)=h^{(j_1)}\cdots h^{(j_s)}$.  Indeed, for $s=0$ this is
the definition of $\psi$.  Assume (1) for $s$.  Each $D^\nu g$ is $C^\infty$ on
$\Omega$, and by the chain rule, for any $C^1$ function $F$ on $\Omega$,
\[
  \frac{\mathrm{d}}{\mathrm{d}t}F(w+th)=\mathrm{d}F(w+th)[h]
  =h^{(0)}\,\mathrm{d}F(w+th)[e_0]+h^{(1)}\,\mathrm{d}F(w+th)[e_1]
  =\sum_{j\in\{0,1\}}h^{(j)}(\partial_jF)(w+th),
\]
using the $\mathbb{R}$-linearity of $\mathrm{d}F(w+th)$ in the direction.
Applying this to each $F=D^\nu g$ in (1) and using
$\partial_j D^\nu=D^{(j,\nu)}$ --- which is how \noderef{PartialWord} composes,
the word being extended on the left --- gives (1) for $s+1$.  This proves the
claim.

\emph{Step 2: all partials of order at least $m+1$ vanish on $U$.}  We show by
induction on $s\ge m+1$ that $D^\nu g=0$ on $U$ for every $|\nu|=s$.  For
$s=m+1$ this is the hypothesis.  Assume it for $s$ and let $|\nu|=s+1$.  Write
$\nu=(j,\nu')$ with $|\nu'|=s$, so that $D^\nu g=\partial_j(D^{\nu'}g)$.  By the
inductive hypothesis $D^{\nu'}g$ vanishes identically on the open set $U$, so
its real Fréchet derivative vanishes at every point of $U$, and therefore
$D^\nu g=0$ on $U$.

\emph{Step 3: $g$ is a polynomial on each ball contained in $U$.}  Let $z_0\in U$
and let $r>0$ be such that the open ball $B=B(z_0,r)$ is contained in $U$.  Fix
$z\in B$ and put $h=z-z_0$ and $\psi(t)=g(z_0+th)$; the segment
$\{z_0+th:0\le t\le1\}$ lies in the convex set $B\subseteq U$.  By Step 1,
$\psi$ is $C^\infty$ on a neighbourhood of $[0,1]$, and by (1) with $s=m+1$
together with the hypothesis,
\[
  \psi^{(m+1)}(t)=\sum_{|\nu|=m+1}c_\nu(h)\,D^\nu g(z_0+th)=0
  \qquad(0\le t\le1).
\]
Taylor's formula with integral remainder for the one-variable $C^{m+1}$
function $\psi$ on $[0,1]$ therefore gives
\[
  g(z)=\psi(1)=\sum_{s=0}^{m}\frac{\psi^{(s)}(0)}{s!}
    +\int_0^1\frac{(1-t)^m}{m!}\psi^{(m+1)}(t)\,\mathrm{d}t
  =\sum_{s=0}^{m}\frac{1}{s!}\sum_{|\nu|=s}c_\nu(z-z_0)\,D^\nu g(z_0),
\]
where the second equality uses (1) at $t=0$.  Each $c_\nu(z-z_0)$ is a product
of $s$ factors, each equal to $\operatorname{Re}(z-z_0)$ or
$\operatorname{Im}(z-z_0)$.  Writing
\[
  \operatorname{Re}(z-z_0)=\tfrac12\bigl((z-z_0)+\overline{(z-z_0)}\bigr),
  \qquad
  \operatorname{Im}(z-z_0)=\tfrac1{2i}\bigl((z-z_0)-\overline{(z-z_0)}\bigr),
\]
and expanding, and then expanding the powers of $z-z_0$ and
$\bar z-\bar z_0$ by the binomial theorem, we obtain coefficients
$a_{ij}(z_0)\in\mathbb{C}$, vanishing unless $i+j\le m$, with
\begin{equation}\tag{2}
  g(z)=\sum_{i+j\le m}a_{ij}(z_0)\,z^i\bar z^{\,j}\qquad(z\in B(z_0,r)).
\end{equation}

\emph{Step 4: rigidity.}  We claim that if coefficients $b_{ij}$, $i+j\le m$,
satisfy $\sum_{i+j\le m}b_{ij}z^i\bar z^{\,j}=0$ for every $z$ in some nonempty
open set $O\subseteq\mathbb{C}$, then every $b_{ij}=0$.

Let $\mathcal{P}_m$ denote the complex vector space of functions
$\mathbb{R}^2\to\mathbb{C}$ of the form $\sum_{a+b\le m}\gamma_{ab}x^ay^b$.  Its
dimension is the number $N$ of pairs $(a,b)$ with $a+b\le m$, because the
monomials $x^ay^b$ are linearly independent as functions (a nonzero polynomial
in two real variables is not the zero function).  Under $z=x+iy$ each
$z^i\bar z^{\,j}$ with $i+j\le m$ lies in $\mathcal{P}_m$, and conversely each
$x^ay^b$ with $a+b\le m$ is, by the two displays of Step 3, a linear
combination of the $z^i\bar z^{\,j}$ with $i+j\le m$; so the $N$ functions
$z^i\bar z^{\,j}$, $i+j\le m$, span the $N$-dimensional space $\mathcal{P}_m$
and hence form a basis of it.

Now $Q=\sum_{i+j\le m}b_{ij}z^i\bar z^{\,j}$ is an element of $\mathcal{P}_m$
vanishing on the nonempty open set $O$.  Writing $Q=\sum_{a+b\le m}
\gamma_{ab}x^ay^b$ and fixing a point $(x_0,y_0)\in O$, all partial derivatives
$\partial_x^a\partial_y^bQ$ vanish at $(x_0,y_0)$, since $Q$ vanishes on a
neighbourhood of it; and those derivatives at $(x_0,y_0)$ determine all the
$\gamma_{ab}$ (by downward induction on $a+b$, the derivative of order $(a,b)$
at $(x_0,y_0)$ equals $a!\,b!\,\gamma_{ab}$ plus a combination of
$\gamma_{a'b'}$ with $a'+b'>a+b$).  Hence every $\gamma_{ab}=0$, so $Q=0$ in
$\mathcal{P}_m$, and since the $z^i\bar z^{\,j}$ form a basis, every
$b_{ij}=0$.  This proves the claim.

\emph{Step 5: the coefficients do not depend on the base point.}  For $z_0\in U$
choose $r(z_0)>0$ with $B(z_0,r(z_0))\subseteq U$ and let
$\Phi(z_0)=\bigl(a_{ij}(z_0)\bigr)_{i+j\le m}$ be the coefficient family of (2);
it is well defined, because by Step 4 two families representing $g$ on the same
ball coincide.  We claim $\Phi$ is locally constant.  Let $z_0\in U$ and
$z_1\in B(z_0,r(z_0))$.  The set
$B(z_0,r(z_0))\cap B(z_1,r(z_1))$ is open and contains $z_1$, hence is
nonempty, and on it both representations (2) hold, so
\[
  \sum_{i+j\le m}\bigl(a_{ij}(z_0)-a_{ij}(z_1)\bigr)z^i\bar z^{\,j}=0
\]
there.  By Step 4, $a_{ij}(z_0)=a_{ij}(z_1)$ for all $i,j$.  Thus $\Phi$ is
constant on $B(z_0,r(z_0))$, and so locally constant on $U$.

The set $U$ is connected and nonempty, so a locally constant function on it is
constant: $\Phi\equiv(a_{ij})$ for a single family $a_{ij}$, $i+j\le m$.
Extend it by $a_{ij}=0$ for $i+j>m$.  Finally, for $z\in U$ apply (2) with
$z_0=z$, which is legitimate because $z\in B(z,r(z))$:
\[
  g(z)=\sum_{i+j\le m}a_{ij}\,z^i\bar z^{\,j}
      =\sum_{i=0}^{m}\sum_{j=0}^{m}a_{ij}\,z^i\bar z^{\,j},
\]
the second equality because the added terms have $i+j>m$ and so vanishing
coefficients.  This is the assertion.
\end{proof}
